\section{Introduction}

The United States military is, by some measures, the largest geographically mobile institution in the country. Active-duty personnel are assigned to installations by orders they cannot refuse; they live where the military places them, often far from the states and communities they call home. Their dependents follow. When the decennial Census enumerates the population, it must make a choice: count military personnel where they are physically located, or count them where their home-of-record indicates they belong.

The answer is neither simple nor uniform. For domestic personnel stationed in the United States, the Census counts them at their installation residence---the barracks, base housing, or off-base apartment where they live while stationed. For overseas personnel (deployed abroad, stationed at foreign bases, or on ships at sea), the Census uses a different rule: it counts them at their home-of-record state for the purpose of congressional apportionment, then allocates them to their home-state at the state level for apportionment purposes, but does not allocate them to specific tracts for redistricting purposes.

The result is a two-tier treatment with distinct redistricting consequences:

\begin{enumerate}
    \item \textbf{Domestically stationed military} are counted at their physical location, inflating the population of installation-adjacent census tracts and congressional districts.
    \item \textbf{Overseas military and dependents} are counted at the state level but do not add to any specific census tract, meaning they do not directly inflate any specific congressional district.
\end{enumerate}

Both effects have measurable redistricting consequences. This paper focuses primarily on the first---domestic installation concentrations---because it produces the tract-level distortions that propagate through algorithmic redistricting. We address the overseas population as context and for the purpose of explaining BISECT's data handling.

\subsection{Scale of the Population}

As of the 2020 Census, the military and dependent population relevant to redistricting includes:
\begin{itemize}
    \item Approximately 1.4 million active-duty personnel stationed domestically
    \item Approximately 700,000 dependents of domestically stationed active-duty personnel
    \item Approximately 600,000 overseas active-duty personnel and their dependents (counted at home-of-record state, not at tract level)
    \item Approximately 800,000 National Guard and Reserve personnel (counted at their civilian residence, not installation; not discussed further here)
\end{itemize}

The approximately 2.1 million domestically stationed military and dependents are distributed across hundreds of installations in every state, but with extreme concentration: the 20 largest installations account for approximately 60\% of the total.

\subsection{Research Questions}

\begin{enumerate}
    \item How does the Census Bureau's military population counting framework affect census-tract-level population weights, and therefore BISECT's graph partitioning?
    \item Which states and districts are most affected by military population concentration?
    \item What is the partisan direction of military population counting, and how does it interact with algorithmic redistricting's partisan neutrality?
    \item How should BISECT handle military populations in sensitivity analyses?
\end{enumerate}
